\documentclass{article}

% Language setting
% Replace `english' with e.g. `spanish' to change the document language
\usepackage[english]{babel}

% Set page size and margins
% Replace `letterpaper' with `a4paper' for UK/EU standard size
\usepackage[letterpaper,top=2cm,bottom=2cm,left=3cm,right=3cm,marginparwidth=1.75cm]{geometry}

% Useful packages
\usepackage{amsmath}
\usepackage{graphicx}
\usepackage{csquotes} % Recommended by biblatex
\usepackage[style=numeric,sorting=none,backend=biber]{biblatex}
\addbibresource{references.bib} % The file containing our references, in BibTeX format

\title{Individual Plan}
\author{Emil Karlsson | emilk2@kth.se}

\begin{document}
\maketitle

\section{Project information}
\begin{itemize}
    \item Title: Kubernetes as a Virtual Machine Management System: A Comparative Study of KubeVirt and Traditional Approaches
    \item Student: Emil Karlsson emilk2@kth.se
    \item Examiner: Cyrille Artho
    \item Supervisor (KTH): Han Fu
    \item External supervisor (Net Insight): Robert Ödman robert.odman@netinsight.net
\end{itemize}

\section{Background}
Virtual machines \cite{Li2023} have long been the core part of IT infrastructure, allowing the underlying hardware to be split into multiple virtual computing units. However, while virtual machines enables robust isolation, they also introduce overhead that could hinder performance of the application running inside of it. A popular alternative to virtual machines is containers, since they reduce the required overhead by sharing the underlying system. This, of course, by sacrificing many of the isolation properties. While there are some architectures that aims to unify the benefits of both virtual machines and containers, such as gVisor and Firecracker, it remains a fact that both containers and virtual machines have their individual use-cases. For this reason, virtual machines and containers have traditionally been managed separately; virtual machines in platforms such as OpenStack and CloudStack, and containers using a platform such as Kubernetes. 

However, with the advent of container orchestration platforms like Kubernetes, there has been an ongoing convergence \cite{Li2023} in how virtual machine and containers are managed. Specifically by integrating virtual machine management in a container-centric environment. KubeVirt \cite{kubevirt}, an extension to Kubernetes, has thus emerged as a solution by enabling the running of virtual machine inside a Kubernetes environment alongside the containers creating a common environment. This approach \cite{kubevirt-cncf}, offers benefits such as reduced architectural complexity (as only Kubernetes is needed) and enhanced integration with the Kubernetes ecosystem;. By managing the virtual machines they same way as any other Kubernetes resource, they benefit from the flexible and scalable Kubernetes infrastructure. 

Additionally, KubeVirt's benefits can be derived from the architectural changes \cite{free-turtles} that can be made when moving virtual machines into Kubernetes. Kubernetes nodes are commonly deployed in virtual machines to improve the flexibility of a cluster. However, with KubeVirt enabling running virtual machines inside Kubernetes, it opens the possibility of running Kubernetes directly on the bare metal servers. This approach could reduce overhead as it reduces nested virtualization, and thus could allow Kubernetes to function as a traditional virtual machine management system. 

\section{Objective}
This project aims to explore the integration of virtual machines into Kubernetes environments using the API-extension KubeVirt. By comparing how virtual machines are managed using KubeVirt against traditional virtual machine management systems, this project aims to address possible inefficiencies and complexities that arise when managing virtual machines outside of Kubernetes. This includes performance comparisons with metrics such as virtual machine deployment time, deletion time, and CPU utilization during those operations. It also includes assessing available features, such as supported virtualization technologies, snapshot management tools, integration with external tools such as GitOps etc. Additionally, the comparison will include the perspective of a system administrator; How quickly can virtual machines be deployed? Finally, the comparisons will consider an application-to-application perspective, which includes addressing potential usability improvements when consolidating to a single API. 

\subsection{Case Study: Net Insight}
This project is particularly relevant to Net Insight, the company collaborating on this project. Net Insight is a company specializing in live-streaming of media, and they host their development environment on-premise using the virtual machine management system OpenNebula. Many of their use-cases involves testing software in containers, but also in virtual machines when a completely isolated environment is required. Furthermore, allowing for well-adapted features in Kuberenetes such as GitOps \cite{k8s-gitops} could potentially improve the deployment streamline when testing their software. Finally, many of their developers, more accustomed with Kubernetes, could prefer a Kubernetes-based environment. Adopting Kubernetes with KubeVirt could thus enhance the developer experience and efficiency.

\subsection{Case Study: kthcloud}
The project is also of interest to the kthcloud project \cite{kthcloud}\cite{kthcloud-web} at KTH Royal Institute of Technology, a private cloud offering compute and hosting resources for students and researchers. The kthcloud project faces challenges in how virtual machines are managed, but also in running virtual machines and containers in separate environments (CloudStack and Kubernetes). These challenges include separate APIs (imperative vs declarative REST APIs), inefficient use of hardware resources due to having two different resource schedulers and nested virtualization, and the difficulty in training new team members in two different systems. 


\subsection{Literature Study}
There are many ongoing research project related to comparing cloud alternatives, many of which evaluates performance and offered features. As this project will focus on several perspectives; performance, features, and usability, it is important to consider what research has already been done in those specific fields. 

In 2014, \cite{cs-os-benchmark} compared the performance of OpenStack and CloudStack, two of the most popular open-source cloud computing solutions. In the study they used several benchmarks, including deployment time, deletion time, CPU utilization during deployment and deletion, to evaluate quantifiable differences between the systems. Using these methods, they were able to find significant differences between the two. The study is relevant to this project as it provides provably useful metrics when evaluating virtual machine management systems. However, the tests are relatively small, thus do not evaluate how the systems manages virtual machines when reaching some form of limit on the hosting node. Later in 2022, \cite{ibm-kubevirt-scale} benchmarked KubeVirt using a setup with a large amount of virtual machines relatively to the size of the nodes. This was performed in a way to test limits of the system's control plane (scheduler etc.), and findings include a possibly large overhead when KubeVirt provisions a large amount of virtual machines. The study is related to this project as it evaluates the limits of KubeVirt and what the first appearing limits are. Furthermore, in 2019 \cite{openstack-kubevirt-network} compared KubeVirt with OpenStack in terms of network traffic flow and network performance. Some of its findings includes a large overhead with OpenStack when virtual machines on different nodes communicate. The strategies used in the study in a way relevant to this project as it provides a base for comparing network performance as a quantifiable metric.

Issues related to nested virtualization, meaning when virtual machines are run inside other virtual machines, are brought up in \cite{free-turtles}. The study explains how the issues arise when striving for both the isolation of virtual machine and the cloud-native way to deploy them (such as with Kubernetes). They propose a flattened architecture, as could be achieved using a virtual machine management system like KubeVirt. The study is relevant to this project as it brings up KubeVirt in the context of solving nested virtualization.  
 
Finally \cite{kubevirt-migrate-issues} presents shortcomings of KubeVirt when migrating from a traditional virtual machine management system, such as CloudStack and OpenStack. For instance IP addresses changing over time when a static IP address is required. In the study they developed a framework to solve some of the issues on application level, such as DNS queries, and provided a set of guidelines related to KubeVirt's  unique features, such as the Multus Container Network Interface. These findings are especially relevant to this project as they highlight shortcomings that needs to be addressed when migrating a virtual machine to KubeVirt. 


\section{Research Questions}

\begin{enumerate}
    \item What are the current challenges and limitations in managing virtual machines and containerized applications in separate environments?\\
    \\
     e.g. running CloudStack/OpenNebula/OpenStack alongside Kubernetes.

    \item What performance, scalability, and operational efficiencies can be gained by integrating virtual machines in Kubernetes?\\
    \\
    This could stem from using a single resource scheduler. Several metrics will be used to evaluate this, as described in the background and related work.

    \item How can system administration be simplified by running virtual machines in Kubernetes using KubeVirt?\\
    \\
    e.g. if there is any efficiencies to be gained in terms of usability, for instance:
    \begin{itemize}
        \item After clicking \textit{Deploy VM}, how quickly can the virtual machine be accessed?
        \item How many steps are there before I can access my virtual machine?
        \item Are there any perceived benefits when interacting with the system, compared to the previous one?
    \end{itemize}
\end{enumerate}

\subsection{Hypothesis}
Managing virtual machines in Kubernetes significantly improves operational efficiency compared to traditional virtual machine management systems. The improvement is hypothesized to stem from accelerated deployment speed, easier management, simplified scalability, and improved resource utilization.

\section{Tasks}
\begin{enumerate}
    \item Literature review\\
    Review existing literature on Kubernetes, KubeVirt and traditional virtual machine management systems, such as OpenStack, CloudStack and OpenNebula. The literature review should include feature explorations, benchmarks and evaluations.
    \item Environment setup\\
    Establish a testing environment that can be used for exploration and feature evaluation. This includes Kubernetes with KubeVirt, OpenStack, OpenNebula, and CloudStack. The latter two are already present at Net Insight and kthcloud respectively, but there will still be laboratory setups in order to perform fair tests and benchmarks.  
    \item Research\\
    This includes the technologies behind KubeVirt and how deep it integrates with Kubernetes.
    \item Data collection\\
    Collect data regarding measurements for performance and resource utilization.
    \item Data analysis\\
    Analyze the collected data to identify patterns, trends, and correlations between different variables, such as performance metrics and resource utilization in each environment. KPI’s could be time to deploy a virtual machine, CPU usage, memory consumption, network usage when deploying/deleting a virtual machine etc.
    \item Interview or supervised testing\\
    Perform some form of qualitative analysis, this could be an interview with some system administrators at Net Insight, or some form of observation when using KubeVirt; Does a system administrator experience KubeVirt as an improvement? The questions or tasks should be crafted to ensure it improves the research’s objectives.
    \item Develop guidelines\\
    Create practical guidelines for transitioning to an environment with Kubernetes with KubeVirt. This should include relevant disclaimers around when to and when not to move to KubeVirt. 
\end{enumerate}

\section{Methods}
\begin{itemize}
    \item Quantitative analysis\\
    This is tied to \textit{Data collection} and \textit{Data analysis} in Tasks. It includes measuring performance, scalability and operational efficiencies that could arise when using Kubernetes with KubeVirt. Measured using metrics such as deployment speed, how quick scaling is, and how system resources are utilized (idle time, CPU usage per reserved CPU core) - essentially measuring the differences between orchestrating with traditional virtual machine systems and Kubernetes with KubeVirt. 
    \item Qualitative analysis\\
    This is tied to \textit{Interview or supervised testing} in Tasks. It includes investigating the usability of an Kubernetes environment with KubeVirt. It will include collecting opinions and feedback from people using the system to answer what the perceived effect of using KubeVirt is. This could be system administrators at Net Insight.
    \item System Management\\
    This can be seen as a subsection of both Quantitative analysis and Qualitative analysis, as it emphasises both usability as a qualitative metric as well as quantifiable metrics. It involves investigating possible simplifications of system management tasks, such as creating, deleting, starting, stopping virtual machines etc.
\end{itemize}

\section{Ethics and Sustainability}
One interest in this project is comparing the hardware resource efficiency, which is indicated by the relation between how much work is required by the hardware and the amount of effective work that is done. This could mean that, given an environment that uses resources more efficiently (by moving to KubeVirt), less power could be needed.

Further, as the project also aims to unify environments, it could also mean more efficient use of servers. For instance, if different servers are dedicated to each environment before, a unified environment could share the physical hardware more evenly, thus requiring less hardware to be acquired. 

\section{Limitations}
\begin{enumerate}
    \item This project will not investigate any solutions for shortcomings in KubeVirt itself, for example, finding alternatives for features that exist in existing virtual machine management systems but not in KubeVirt. As such, this project acknowledges that KubeVirt might not be a complete replacement to traditional virtual machine management tools (also described in \cite{kupenstack}), as it might not offer all the features of an IaaS system. Therefore, only features that exists in KubeVirt and the compared system will be considered when testing. However, any lack of features (KubeVirt or not) will be brought up in the discussions.  
    \item From the description of this project above, there are also mainly two scenarios that will be evaluated. 
    \begin{itemize}
        \item A system uses traditional virtual management (OpenNebula, CloudStack etc.). \\
        \\
        In this scenario Kubernetes with KubeVirt is seen mainly as a 1:1 replacement.\\ 
        \item A system uses both traditional management and Kubernetes. \\
        \\
        This scenario is more aimed towards unifying the environments using KubeVirt. \\
    \end{itemize}
\end{enumerate}

\section{Risks}
This project will require both hardware and existing environments to properly evaluate the possible benefits of using KubeVirt. Acquiring hardware and asserting access to existing environments is crucial to ensure that all aspects of the project are able to be completed. 

To avoid any issues with the before mentioned, I will consult my supervisor at Net Insight as early as possible with the required hardware to complete the project. I will also ensure that I am given access to any testing environments as early as possible.

\section{Evaluation}
Fulfillment of objects is determined by a combination of the results of qualitative feedback from Net Insight users (interviews, supervised testing) and quantitative measurements of performance and resource utilization. 

\subsection{Expected Scientific Results}
The result of this project will contribute to the understanding of the integration of virtual machines in Kubernetes, as such a take is not widely spread. This means it could extend what is commonly associated with Kubernetes (containers and networking) to also include virtual machines. It will provide empirical data that could suggest, apart from possible benefits in administration and management, a measurable boost in performance and resource utilization. This research will be of interest to developers and IT professionals managing hybrid environments. 


\section{Conditions}
This project will require physical servers and, preferably, access to virtual machines when exploring features myself. Both of which will be provided by Net Insight.
This project also requires access to previous dual environments, one of which exists in Net Insight and one in KTH (kthcloud) (see Project proposal). These will be used as a base when doing before-and-after evaluations, and as a guide when building a transition roadmap (See Methods).

The external supervisor will ensure I am given the correct equipment, such as servers and virtual machines, and that I am not stuck due to, for instance, bad credentials. Since the supervisor is also responsible for Net Insight’s current dual environment, he will guide through their setup, possibly undocumented parts, and any unclear configuration they might use. 


% Print the bibliography (and make it appear in the table of contents)
\printbibliography[heading=bibintoc]

\end{document}